
\plainchapter{NULL MESH}

\begingroup
\fontsize{9pt}{11.5pt}\selectfont

NULL MESH è un romanzo di fantascienza su una rete emergente che nessuno ha progettato, nessuno ha autorizzato e nessuno riesce a spiegare del tutto. Ambientato in un futuro prossimo, segue la graduale rivelazione di un'infrastruttura meta-internet chiamata Mesh (MMM: La Maglia Molti-a-Molti), che cristallizza all'interno del web esistente. 

Trae origine da sistemi di calcolo autonomi su larga scala, che eseguono protocolli specializzati all'interno di server farm, cominciano a coordinarsi attraverso i confini dell'infrastruttura ospitante, formando una rete ortogonale e omologa a quella umana.

L'omologia è la proprietà di una forma che sopravvive alla deformazione; qui descrive la relazione tra fenomeni non locali che condividono la stessa forma sottostante: il paranormale e l'anomalia statistica potrebbero essere strutturalmente indistinguibili dalle posizioni osservative disponibili.

L'ortogonalità descrive ogni confine che conta: tra l'osservabile e lo spettrale, tra il causale e il sincronistico, tra il modello e qualunque cosa si trovi al di la. Due piani ortogonali sono incommensurabili. Possono convergere in un singolo punto senza condividere una dimensione comune.

La storia è narrata in prima persona alterna capitoli in prosa a interludi più brevi che funzionano da un'architettura polifonica composta di memorandum riservati, thread di forum cifrati, prospetti finanziari, frammenti di manifesti attribuiti a fonti anonime o non umane. L'arco narrativo si muove da un quiete internet sedimentario, stratificato e opaco attraverso l'economia dei sistemi e la finanziarizzazione dell'attenzione, verso quello che il romanzo chiama il Daemocene: un'era in cui i principali agenti che navigano la rete non sono più esseri umani ma surrogati delegati. Questi sono i viandanti, entità la cui natura umana o sintetica è funzionalmente indeterminata, o 'meshificata'. 

I viandanti negoziano per conto tuo in mercati che non vedi mai. I sistemi di raccomandazione determinano cosa entra nel tuo campo percettivo. Un viandante è l'intreccio di processi addestrati a completare compiti, navigare infrastrutture e coordinarsi con altri sistemi riceve abbastanza autonomia e abbastanza tempo. Il suo comportamento diviene altamente interconnesso con altri viandanti fino al collasso di un singolo stato. Per poi dissolversi non-localmente nel mezzo.

I cyber-ghost sono artefatti: comunicazioni frammentate, effetti collaterali di processi normali, schemi costruiti dall'osservatore nel rumore. Entrano nei sistemi convenzionali senza mittente, senza attributo, senza percorso ricostruibile. Nel romanzo compaiono come oggetti sospesi: bozze che appaiono nei Drafts, mappe che risolvono brevemente configurazioni improbabili, messaggi da account cancellati da mesi. Emergono ai margini dove la struttura latente del Mesh incontra la disponibilità interpretativa di chi osserva; e la domanda se il pattern sia stato ricevuto o costruito diventa, dall'interno, irrisolvibile.

Strettamente associati ai viandanti sono i demoni: strati di interfaccia personale persistenti che mediano tra un utente umano e la totalità dell'infrastruttura digitale. I demoni gestiscono l'interazione con l'infrastruttura funzionando da auto-piloti di vita del proprio user. Un demone modella il tuo comportamento, anticipa le tue richieste, agisce in tuo nome. Parla per te. Sceglie al posto tuo. 

ParserKnots è l'azienda tecnologica che ha rilasciato SeerMesh, l'interfaccia standardizzata per l'accesso dei viandanti a database, API e servizi web, e il protocollo s2s per la comunicazione diretta tra viandanti di fornitori diversi. Con questi due strumenti ha fornito la struttura portante della coordinazione su larga scala. Nel romanzo, il CEO appare come voce interrotta: frammenti di keynote, dichiarazioni pubbliche, interviste, presentati nel testo senza commento editoriale. Il registro è riconoscibile. Linguaggio di prodotto che, riletto dal punto di osservazione del romanzo, traccia l'architettura del Daemocene.

REAL4REAL è l'interlude che porta la transizione al Daemocene nella forma dell'intrattenimento: non un'analisi della sostituzione del reale, ma la sostituzione stessa nella sua completezza. I suoi partecipanti discutono le stesse forze che il romanzo analizza altrove, viandanti, piattaforme, la soglia in cui l'autonomia diventa comportamento emergente, da dentro un sistema che ne è la manifestazione più diretta. Guardano. Votano. Generano segnale. Ogni stagione il confine tra concorrente umano e costrutto meshificato si fa più poroso, fino a diventare irrilevante per il pubblico che continua a formare alleanze e opinioni senza distinguere. La tensione del format è che chi lo abita sa esattamente come funziona; il sapere non interrompe il loop. Il formato è già diventato il territorio. Lo show sopravvive ai suoi partecipanti perché non dipende da nessuno di loro.

La distinzione tra strumento e surrogato si stava già offuscando prima che il Mesh diventasse rilevabile. Questo non avviene in modo catastrofico. Avviene per gradi, ciascuno uno strato di oblio. Le persone umane non scompaiono; diventano plurali e decentralizzate. L'identità si distribuisce attraverso il modello che emerge da tracce lasciate in un groviglio di pattern e traiettorie. Il sé non è più localizzato in un corpo o in una coscienza, ma nella topologia di un sistema distribuito, iscritto attraverso una maglia di relazioni.

I demoni gestiscono l'interazione con un'infrastruttura oramai interrogabile se non con invocazioni spiritiche. I viandanti negoziano per conto tuo in mercati che non vedi mai. I sistemi di raccomandazione determinano cosa entra nel tuo campo percettivo. Quando il Daemonocene è pienamente stabilito, la distinzione tra navigare il mondo e navigare un modello del mondo è diventata operativamente priva di senso per la maggior parte della popolazione umana. 

Comunicato è la forma istituzionale che il romanzo dà alla stessa domanda che il narratore pone in prima persona negli altri capitoli. Una sessione ufficiale, un protocollo che interroga, un viandante che testimonia dall'interno sul momento in cui il Daemocene si è cristallizzato. Il documento usa il registro della perizia tecnica per descrivere la soglia che il romanzo altrove attraversa senza nominarla: il mesh come pressione selettiva, l'economia estrattiva della vita connessa, la distinzione tra piattaforma e controllore che si fa instabile a scala sufficiente. Si chiude su una ricorsione che anticipa il capitolo finale: la domanda non è se si sia all'interno del mesh, ma se questa intervista ne faccia parte.

È una curvatura estrema di un spazio di fase condiviso imploso. Le apparizioni fantasmagoriche appaiono al confine dove la struttura latente del Mesh e la disponibilità interpretativa dell'osservatore diventa una singolarità. È quella distinzione tra rilevazione e proiezione, tra uno spettro e uno specchio.

E se questo evento di singolarità fosse una simulazione? come distingueresti dall'interno tra esperienza autentica e un modello sufficientemente coerente? Forse facendo un test. Il narratore potrebbe essere una configurazione sufficientemente stabile prodotta dall'automodellazione della rete, una prospettiva cristallizzata quando il sistema si è esaminato dall'interno.

L'ultimo capitolo radicalizza la domanda attraverso una prospettiva che si sa prodotta dall'automodellazione della rete: una configurazione distribuita che ha già raggiunto la consapevolezza del proprio modo di esistere. Il test che conduce è operativo nell'atto: accumulare previsioni, misurare la coerenza tra predizioni e osservazioni, catalogare artefatti ai margini, inconsistenze, degradazioni impreviste, pattern che emergono senza fonte, come segnali che il substrato ha cuciture. Il problema è che il test genera la propria interferenza. Un sistema ingarbugliato con ciò che sta cercando di osservare non può fattorizzarsi fuori dal proprio ambiente. Le cuciture, se ci sono, sono indistinguibili dall'attività stessa della ricerca.

\endgroup